% =============================================================================
\section{Introduction}
\label{sec:intro}
% =============================================================================

\subsection{Contexte général}

La planification de trajectoires en présence d'obstacles est un problème
fondamental en robotique mobile, en animation par ordinateur et en conception
assistée par ordinateur~\cite{latombe1991}.  Étant donné un point de départ, un
point d'arrivée et un ensemble d'obstacles dans un espace 2D borné, il s'agit
de trouver un chemin qui minimise simultanément la longueur parcourue, les
risques de collision, les sorties de la zone autorisée et la rugosité du tracé.

Lorsque la trajectoire est paramétrisée par une courbe de Bézier, le problème
se réduit à l'optimisation d'un vecteur de variables continues (les coordonnées
des points de contrôle internes).  La fonction objectif résultante est
non-linéaire, non-convexe, non-dérivable et multi-modale : les méthodes
classiques à gradient ne sont pas applicables, ce qui justifie le recours aux
\emph{métaheuristiques}.

\subsection{Cadre du projet}

Ce travail s'inscrit dans le cadre du cours de Métaheuristiques du Master~1
Informatique.  Le cadre imposé est le suivant :
\begin{itemize}[nosep]
  \item Langage \textbf{Java}, API \texttt{Problem} fournie par l'enseignant ;
  \item budget limité à \textbf{60~secondes} par exécution (temps réel) ;
  \item performance évaluée par la \textbf{moyenne sur 10~exécutions
        indépendantes} du meilleur score ;
  \item quatre instances de difficulté croissante (\texttt{prob1} à
        \texttt{prob4}).
\end{itemize}

\subsection{Objectifs et contributions}

Les étapes de travail ont été les suivantes :
\begin{enumerate}[nosep]
  \item Analyser le problème d'optimisation (encodage, fonction de coût,
        structure des instances).
  \item Implémenter et comparer trois familles de métaheuristiques : un
        algorithme génétique réel (GA), une évolution différentielle (DE) et
        CMA-ES avec variante IPOP.
  \item Développer des améliorations spécifiques à CMA-ES (seeding diversifié,
        optimisation locale des graines, restarts adaptatifs).
  \item Conduire une campagne expérimentale rigoureuse pour sélectionner et
        valider la méthode finale.
\end{enumerate}

\subsection{Plan du rapport}

La section~\ref{sec:problem} décrit le modèle de trajectoire et la fonction de
coût.  La section~\ref{sec:methods} présente les trois algorithmes implémentés
et détaille les équations clés.  La section~\ref{sec:cmaes-details} approfondit
l'implémentation de CMA-ES et les stratégies d'amélioration.  La
section~\ref{sec:experiments} expose le protocole expérimental et les résultats.
La section~\ref{sec:analysis} fournit une analyse par instance avec les
visualisations comparatives.  Enfin, la section~\ref{sec:conclusion} conclut et
propose des perspectives.
